Each clause is the imported structure read under the premise. (1) is the inaccessibility of Premise \claimref{one-relation.recorded-sectors-superselected-premise} (the paper's Remark on recording) applied to a subject constituted by records in one sector; the content is that the supervenience (\claimref{one-relation.supervenience-stated-so-it-can-be}) premise turns an operational statement about correlations into a statement about what a self can contain. (2) is the irreversibility of Premise \claimref{one-relation.recording-irreversible-stated-so-it-can} on the modular period of Fact \claimref{one-relation.imported-single-mode-thermality-modular-period}, together with the flag identification of Fact \claimref{one-relation.imported-flag-arrow-superselection}, under Hypothesis (S) (\claimref{standard-model.hyp-s-slot-statistics}) as that Fact records: the writing direction is the flag's orientation, and under the premise the felt direction is the writing direction. (3) is the self-location argument: a subject individuated by a record chain has no counterpart in a sibling chain to be identified with, so the indexical question has no referent capable of being surprised by its answer. The proposition proves nothing about the hard problem and says so.
